\documentclass[preprint,12pt,authoryear]{elsarticle}
\usepackage{amssymb} % better math symbols
\usepackage{amsmath} % better math environments
\usepackage{lineno}  % 用于在左侧添加行号
\usepackage{natbib}  % 用于引用样式
\usepackage{graphicx}
\usepackage{mathptmx}  % time new roman fonts for pdflatex
% 添加超链接
\usepackage[breaklinks]{hyperref}
\usepackage[hyphenbreaks]{breakurl}

\usepackage{setspace}  % double space
\usepackage{indentfirst}  % first paragraph also indent

\journal{}

\begin{document}

\linenumbers

\begin{frontmatter}
\title{Your Paper Title}

\author{}

\affiliation{organization={},
            addressline={}, 
            city={},
            postcode={}, 
            state={},
            country={}}

\begin{abstract}
Your abstract goes here.
\end{abstract}

% \begin{graphicalabstract}
%   \includegraphics{grabs}
% \end{graphicalabstract}

% \begin{highlights}
% \item Research highlight 1
% \item Research highlight 2
% \end{highlights}

\begin{keyword}
keyword1 \sep keyword2 \sep keyword3
\end{keyword}

\end{frontmatter}

% \doublespacing

\section{Introduction}
\label{sec:introduction}
Your introduction text here \citep{iwahashi_drought_2022}.

\section{Methods}
\label{sec:methods}
Your methods description here.

\begin{equation}
    V = T_e(F)
    \label{eq:3}
\end{equation}

{\raggedright where ...}

{\raggedright\textbf{(1) Centroid-based berry contour completion}}

\section{Results}
\label{sec:results}
Your results here.

\section{Discussion}
\label{sec:discussion}
Your discussion here.

\section{Conclusion}
\label{sec:conclusion}
Your conclusion here.

\section*{Acknowledgement}
Your acknowledgements here.

\bibliographystyle{elsarticle-harv}
\bibliography{references}

\end{document}